\chapter{Methodology}
\label{sec:methodology}
% Focus on what you add to the existing method. Explain what you will do and why (and how). Do not forget to characterize your research design. There should be a sub-section on the evaluation. 
%For DS students, this normally means using manually labelled or ground truth data. For IS students, it is not always needed to have a separate methodology section. You can also integrate the approach with the results in one section. It depends on your type of research what is best fitting.
Write about your methodology here. Focus on your own contribution. Indicate exactly how you will assess your work in terms of evaluation.

% It is possible to use a separate section for the Experimental Setup, which then focuses on all settings used in your experiments. It also possible to address the settings in a sub-section under Methodology. 
\section{Equations}
        We estimate the deformable template parameters~$\theta_t$ and the deformation fields for every data point using maximum likelihood. Letting~$\mathcal{V} = \{\boldsymbol{v}_i\}$ and~$\mathcal{A} = \{a_i\}$,
        %
        \begin{align}
                \hat{\theta_t}, \hat{\mathcal{V}} &= \arg \max_{\theta_t, \mathcal{V}} \log p_{\theta_t}(\mathcal{V} | \mathcal{X},  \mathcal{A}) \nonumber \\
                &= \arg \max_{\theta_t, \mathcal{V}} \log p_{\theta_t}(\mathcal{X} | \mathcal{V}; \mathcal{A}) + \log p(\mathcal{V}),
                \label{eq:logpost}
        \end{align}
        %
        where the first term captures the likelihood of the data and deformations, and the second term controls a prior over the deformation fields.

        \begin{proof}
                Awesome proof.
        \end{proof}

\section{Long equations in two columns}
        Note that equations can fill the width pretty quickly when there are two columns.
        Different mathemtical environments, beyond the basic \verb|\begin{equation}| may be of use. Taking as an example the binomial formula which overflows in Eq. \eqref{binom:eq}:
        \begin{equation}
                \label{binom:eq}
                (1+x)^n = \underbrace{(1+x)\times...\times(1+x)}_{n \text{ times}} = \sum_{k=0}^n \binom{n}{k}x^k = \sum_{k=0}^n\frac{n!}{k!(n-k)!}x^k .
        \end{equation}

        The \verb|\begin{multline}| environment is an easy although crude fix:
        \begin{multline}
                (1+x)^n = \underbrace{(1+x)\times...\times(1+x)}_{n \text{ times}} \\
                = \sum_{k=0}^n \binom{n}{k}x^k \\
                = \sum_{k=0}^n\frac{n!}{k!(n-k)!}x^k .
        \end{multline}
        \verb|\begin{align}| tends to give the most control over the final result:
        \begin{align}
                (1+x)^n &= \underbrace{(1+x)\times...\times(1+x)}_{n \text{ times}} \nonumber\\
                        &= \sum_{k=0}^n \binom{n}{k}x^k \nonumber\\
                        &= \sum_{k=0}^n\frac{n!}{k!(n-k)!}x^k .
        \end{align}


\section{Math styles}
        Different font styles can be used for equations:
        \begin{itemize}
                \item \verb|$a b c A B C 1 2 3$|: $ a b c A B C 1 2 3 $
                \item \verb|$\mathbf{a b c A B C 1 2 3}$|: $ \mathbf{a b c A B C 1 2 3} $
                \item \verb|$\mathfrak{a b c A B C 1 2 3}$|: $ \mathfrak{a b c A B C 1 2 3} $
                \item \verb|$\mathcal{ABC}$|: $ \mathcal{ABC} $
                \item \verb|$\mathbb{ABC}$|: $ \mathbb{ABC} $
                \item \verb|$\mathsf{ABC}$|: $ \mathsf{ABC} $
                \item \verb|$\mathtt{ABC}$|: $ \mathtt{ABC} $
        \end{itemize}

        Text and names in equations should be dealt with the \verb|\text| command, for instance:\\
        \verb|$\mathcal L_{\text{SuperLoss}}$|: $\mathcal L_{\text{SuperLoss}}$ and not $\mathcal L_{SuperLoss}$.